% ! TeX root = ../../main.tex

\section{Lessons for BSP Diagnostic Workflows}\label{sec:amr:concl}

This chapter traced the end-to-end process of developing telemetry-driven placement policies for
AMR codes.
It illustrates the fragility of the BSP paradigm:
a single straggler at any
synchronization point idles all other ranks, and the likelihood of encountering one grows with
scale.
% Diagnosing and mitigating such effects requires sustained visibility across the full stack.
The technical outcome of this study was CPLX, a placement policy that reduced runtime by up to
21.6\% over tuned baselines by exposing the tradeoff between load balance and communication
locality.
But the dominant effort was not placement design---it was the diagnostic process: identifying and
eliminating confounders across the stack, evolving our telemetry and analysis workflow through
several iterations, and repeatedly validating that observed effects were genuine.
The following lessons distill what this process revealed about the requirements for effective BSP
diagnostic workflows, and inform the design of ORCA in the next chapter.

\lesson{1}{Straggler attribution requires fine-grained telemetry}

\noindent Distinct straggler mechanisms are indistinguishable without access to fine-grained
telemetry.
Load imbalance, hardware variability, scheduling effects, and tuning artifacts all manifest as
late arrival at a barrier.
Aggregate timing alone cannot distinguish between them---attribution requires per-rank,
per-timestep visibility into the events leading up to the stall.
Codes with peer-to-peer message patterns add further complexity.
Asynchronous execution introduces scheduling and ordering as additional sources of variability,
and the messages themselves create indirect paths through which delays propagate before surfacing
at synchronization.
Fine-grained runtime telemetry is a prerequisite for disentangling these causes.

\lesson{2}{Attribution requires live runtime context}

\noindent The predominant mechanism for fine-grained collection in BSP is traces captured as
per-rank logs
and analyzed post-hoc.
Traces are effective for general characterization of workload properties, but root-cause
localization often requires additional telemetry and hardware and runtime context that is no longer
available once the
trace has left the execution environment~\cite{Klenk2017:OverviewMPICharacteristics}.
Many performance pathologies can only be diagnosed and resolved in the production environment
where they manifest, which requires that the relevant signals be observable in situ.

\lesson{3}{Mitigations are context-specific and require continuous visibility}

\noindent Runtime unknowns such as tunable parameters encode context-specific decisions and do not admit
universal fixes.
Our placement policy, CPLX, improved runtime by up to 21.6\%, but its optimal tradeoff between
load balance and communication locality varies with workload, cluster, and scale.
Such conditions can be mitigated but not eliminated.
Continuous visibility is needed not only to guide initial parameter selection, but also to detect
regressions as conditions change.

\lesson{4}{Fine-grained telemetry collection must be hypothesis-driven and guided by causal
  structure}

\noindent All static tracing profiles cut off at arbitrary instrumentation depths.
Execution has practically infinite fidelity: every function call can be
traced further.
Tracing deeper exhaustively is both unviable and unnecessary.
Stragglers in BSP, for example, can be identified from collective latency distributions across
ranks.
Only the anomalous ranks warrant deeper inspection, and only the functions contributing to their
delays.
Effective diagnosis must therefore be hypothesis-driven: it follows causal structure downward in
response to observed anomalies, instead of attempting exhaustive collection in advance.

\lesson{5}{Hypotheses require extensible telemetry collection and programmable analytics}

\noindent Fine-grained telemetry must be assembled into views that preserve causal relationships
before hypotheses can be validated.
This places requirements on both the collection and analytics layers.

\begin{itemize}[leftmargin=*]
  \item Collection infrastructure must permit flexible telemetry schemas.
        In our study, while standard tracing tools~\cite{Shend2006:TAU} captured per-task telemetry via
        Kokkos~\cite{Hammo2018:KokkosP}
        instrumentation, they did not capture application-level \emph{block IDs} that we needed to
        connect timing data to meshblocks for placement analyses.
        Telemetry from different runtime sources may require different schemas --- no static schema
        can anticipate every analytical need.

  \item Analytics must support programmable views assembled on demand.
        eBPF demonstrated this paradigm during our investigation: its mechanisms enable selective
        instrumentation of arbitrary points in the execution stack, activated on demand rather
        than compiled in.
        This was decisive for localizing latency spikes when other approaches failed.
        The limitation is that eBPF is node-local---events are indexed by hostname and process
        identifier, probes must be manually activated on each node, and results must be stitched
        across nodes and resolved into ranks.
\end{itemize}

At BSP scales, predominant event trace
formats~\cite{Eschw2012:OTF2,:ChromeJsonTEF,:CaliperPerformanceAnalysis} serve neither timeline
visualization nor analytics.
The size of these traces precludes interactive exploration via timeline-based trace
viewers~\cite{:PerfettoSystemProfiling},
and their structure and formats impose prohibitive ETL (\emph{extract-transform-load}) costs before
dataframe
style queries can be computed.
Columnar and queryable representations~\cite{:Parquet} are the right foundation for managing BSP
telemetry.